\chapter{La carte des plateformes}\label{ch:pl:the-platform-map}

À 07~h~40, le rapport de risque de nuit, le P\&L et le backtest d'un chercheur ne concordent pas sur une même position. Trois équipes passent le
matin à comprendre pourquoi. Le fichier d'opérations sur titres du fournisseur, qui arrive normalement à 19~h~30, est arrivé à 21~h~30 la veille au soir~; le
batch de risque, qui démarre à 21~h~00 selon l'horloge, avait déjà lu l'ancien référentiel des titres (security master) et les positions construites dessus~; le
job de P\&L, qui démarre dès que ses entrées sont prêtes, a lu les nouvelles vingt minutes plus tard. Chaque système a fait ce pour quoi il a été conçu.
Ce qui manquait, c'était une image, tenue sous forme de données et vérifiée par un programme, indiquant quel système lit quoi, auprès de qui et quand. Ce
chapitre dessine cette image pour la société miniature construite tout au long de la série, et en fait un outil (\texttt{firm.platmap})
qui répond aux questions sur lesquelles les trois équipes ont passé leur matinée~: ce qu'atteint un fichier en retard, quel job lit une entrée avant
qu'elle soit complète, et jusqu'à quel retard chaque entrée peut arriver avant que les rapports du matin ne manquent leur échéance.

\section{Les systèmes d'une société de trading}

Une société de trading est une chaîne de systèmes qui transforment les données de marché et les décisions en positions, et les positions en argent qui a
été compté, contrôlé et déclaré. Les livres~1 à 14 de cette série ont construit la plupart des maillons un par un~: le simulateur de bourse
et son flux (livre~10), le feed handler, la passerelle d'ordres et la barrière de risque (livre~13), les moteurs de backtest et le flux de travail
de recherche (livre~7), la bibliothèque de pricing (livre~5), le moteur de risque (livre~6), la plateforme d'apprentissage automatique (livre~12). Ce livre porte sur
les plateformes qui les relient, et il commence par la carte.

\begin{definition}[Flux front-to-back]\label{def:pl:the-platform-map:ftb}
Le \emph{flux front-to-back}\index{flux front-to-back} d'une société de trading est le chemin que suivent une transaction et les données qui la concernent
à travers les systèmes de la société~: données de marché et de référence~; la décision et l'ordre~; l'exécution~; la position et son P\&L~;
le risque qu'elle ajoute~; son enregistrement, sa confirmation, son règlement-livraison et son rapprochement~; et les contrôles et les rapports qui clôturent la journée.
\end{definition}

Le secteur désigne les tronçons de ce parcours d'après les équipes qui les ont historiquement exploitées~: le front, le middle et le back office,
que le livre~16 (chapitre~13) définit comme des unités d'organisation. Ce livre s'intéresse aux systèmes, quelle que soit l'équipe qui les fait tourner.
La \Cref{fig:pl:the-platform-map:layers} montre les dix-huit systèmes de la société miniature, répartis en huit couches, et le sens dans lequel
les données de la journée y circulent.

\begin{omfigure}
\begin{tikzpicture}[font=\scriptsize,>=Stealth,
  sys/.style={draw,rounded corners=2pt,fill=omDef!10,minimum height=0.55cm,inner sep=2pt,align=center},
  ext/.style={sys,fill=gray!12},
  lay/.style={font=\scriptsize\itshape,anchor=east,text=gray!70!black}]
\foreach \y/\name in {0/externe,-1/données,-2/trading,-3/positions,-4/risque,-5/post-marché,-6/contrôle}
  {\node[lay] at (-0.2,\y) {\name};}
\node[ext] (ex) at (1.5,0) {bourses};
\node[ext] (dv) at (4.4,0) {fournisseur de données};
\node[ext] (cb) at (8.8,0) {courtier compensateur};
\node[sys] (tc) at (1.0,-1) {capture\\des ticks};
\node[sys] (ts) at (2.9,-1) {base\\de ticks};
\node[sys] (rd) at (5.0,-1) {données\\de référence};
\node[sys] (ss) at (7.3,-1) {service\\d'instantanés};
\node[sys] (og) at (1.9,-2) {stratégies et passerelle d'ordres};
\node[sys] (ps) at (1.9,-3) {service\\de positions};
\node[sys] (pn) at (4.2,-3) {service\\de P\&L};
\node[sys] (pl) at (4.2,-4) {bibliothèque\\de pricing};
\node[sys] (rg) at (6.8,-4) {grille\\de risque};
\node[sys] (tcap) at (1.9,-5) {saisie des transactions};
\node[sys] (rc) at (8.8,-5) {rapprochement};
\node[sys] (pc) at (1.9,-6) {product control};
\node[sys] (sv) at (3.9,-6) {surveillance};
\node[sys] (rr) at (6.8,-6) {reporting réglementaire};
\node[sys,fill=omMeth!12,minimum height=2.4cm,text width=1.5cm] (rp) at (10.9,-2.7) {plateforme de recherche};
\draw[->] (ex) -- (tc);
\draw[->] (tc) -- (ts);
\draw[->] (dv) -- (rd);
\draw[->] (og) -- (ps);
\draw[->] (ps) -- (pn);
\draw[->] (pl) -- (pn);
\draw[->] (pl) -- (rg);
\draw[->] (ps) -- (tcap);
\draw[->] (cb) -- (rc);
\draw[->] (ts.south) |- (rp.north west |- 0,-1.55);
\draw (rd.south) -- (rd.south |- 0,-1.55);
\draw[->] (ss.south) -- (ss.south |- rg.north);
\end{tikzpicture}

\omcaption{Les systèmes de la société miniature par couche, des plateformes de négociation en haut aux contrôles en bas~; la plateforme de recherche
lit la base de ticks et les données de référence à côté du chemin de trading. Les flèches indiquent le sens principal des données de la journée~; la liste complète
des flux est donnée par les données de la carte (\texttt{pl\_platmap}), et non par le dessin.}\label{fig:pl:the-platform-map:layers}
\end{omfigure}

Lu du haut vers le bas, les couches donnent l'ordre dans lequel naissent les données d'une transaction. Les données de marché et de référence existent avant toute
décision~; la passerelle d'ordres envoie les ordres que la plateforme exécute~; le service de positions conserve le cumul de ces exécutions~; le
service de P\&L la valorise avec les valorisations de la bibliothèque de pricing~; la grille de risque la réévalue sous différents scénarios~; la saisie des transactions l'enregistre et
les systèmes post-trade la confirment, la règlent et la rapprochent~; le product control, la surveillance et le reporting réglementaire contrôlent et déclarent
l'ensemble de ce qui précède. Une plateforme de recherche se tient à l'écart~: elle lit la même base de ticks et les mêmes données de référence pour construire la
prochaine stratégie, et ses sorties reviennent sous forme de paramètres et de modèles que la couche de trading charge (chapitre~16).

\begin{example}[La société miniature en chiffres]\label{ex:pl:the-platform-map:numbers}
La carte de la société miniature (\texttt{pl\_platmap}) compte 18 systèmes, 22 jeux de données et 15 jobs de fin de journée. Sept jeux de données
arrivent de l'extérieur~: les exécutions de la journée, la capture des données de marché de la journée, les cours de clôture officiels, un instantané
de données de marché sur les courbes et les volatilités, les fichiers de données de référence et d'opérations sur titres du fournisseur, et le relevé du courtier compensateur. Chaque
autre jeu de données est produit par exactement un job à partir d'autres.
\end{example}

\section{Flux de données et leurs horloges}

La difficulté d'une plateforme ne tient pas au nombre de systèmes mais au nombre d'horloges. Les données de marché et les exécutions arrivent en continu
pendant la séance~; les cours de clôture officiels arrivent quelques minutes après la clôture~; les fichiers des fournisseurs arrivent le soir, à des heures que le
fournisseur ne garantit pas~; le relevé du courtier compensateur arrive lorsque la chambre de compensation a terminé sa propre journée. Chaque job
de la chaîne de fin de journée fonctionne sur l'une des deux sortes d'horloges~: il démarre dès que ses entrées sont prêtes, ou il démarre à une heure fixe
de l'horloge murale.

\begin{definition}[Carte des flux de données]\label{def:pl:the-platform-map:map}
Une \emph{carte des flux de données}\index{carte des flux de données} est une description, tenue sous forme de données, des systèmes de la société, des jeux de données que chacun détient,
et des jobs qui lisent des jeux de données et en écrivent d'autres, chaque job étant assorti de sa durée et de son horloge (pilotée par les dépendances ou planifiée à une heure
fixe), chaque jeu de données étant assorti de son \omterm{def:pl:the-platform-map:sor}{système de référence} (system of record) et, pour les entrées externes, de son heure d'arrivée.
\end{definition}

Parce qu'elle est constituée de données, une carte peut être vérifiée (chaque jeu de données a un seul propriétaire, les jobs ne forment aucun cycle), parcourue (ce
qu'atteint ce fichier~?) et simulée (à quelle heure chaque rapport sera-t-il prêt ce soir~?). Un diagramme dessiné une fois pour une présentation ne fait aucune de ces opérations, et
est périmé dès la version suivante.

\begin{omfigure}
\begin{tikzpicture}[font=\scriptsize,>=Stealth]
\draw[->] (0,0) -- (12.5,0) node[right]{temps};
\foreach \x/\t in {0/09:30,3.63/16:00,5.58/19:30,6.70/21:30,9.77/03:00,12.0/07:00}
  {\draw (\x,0.08) -- (\x,-0.08) node[below]{\t};}
\fill[omDef!25] (0,0.3) rectangle (3.63,0.75); \node at (1.8,0.525) {flux~: ticks, exécutions};
\fill[omMeth!25] (3.67,0.3) rectangle (10.28,0.75); \node at (7.0,0.525) {batch de fin de journée};
\fill[omThm!20] (6.93,0.95) rectangle (10.28,1.4); \node at (8.6,1.175) {recherche de nuit};
\draw[omThm,thick] (12.0,0.2) -- (12.0,1.6) node[above]{rapports attendus};
\draw[->,thick] (3.63,1.55) node[above]{clôture} -- (3.63,0.8);
\draw[->] (5.58,1.55) node[above,align=center]{fichiers des\\fournisseurs} -- (5.58,0.8);
\draw[->] (6.70,2.45) node[above]{relevé de compensation} -- (6.70,0.8);
\end{tikzpicture}

\omcaption{Les horloges d'une journée de négociation. Ticks et exécutions circulent en flux pendant la séance~; le \omterm{def:pl:the-platform-map:batch}{batch de fin de journée} démarre à la clôture et
attend des fichiers dont il ne contrôle pas les heures d'arrivée~; la recherche de nuit s'exécute après que le batch a produit ses entrées~; les
rapports du matin sont attendus à 07~h~00. Les heures sont les heures nominales de la société miniature.}\label{fig:pl:the-platform-map:clocks}
\end{omfigure}

Le \Cref{lst:pl:the-platform-map:jobs} est la chaîne de fin de journée de la société miniature telle que la carte la consigne~: chaque job avec ses
entrées, sorties et durée en minutes. Le batch de risque accepte un départ fixe facultatif~; tout le reste démarre dès que ses entrées sont disponibles.

\omcode{code/platforms/01-the-platform-map/python/pl_platmap.py}{61}{81}{Les jobs de fin de journée de la société miniature~: entrées,
sorties, durée en minutes et, pour le batch de risque, un départ fixe facultatif.}\label{lst:pl:the-platform-map:jobs}

\section{Systèmes de référence et sources de référence}

Un même nombre existe à plusieurs endroits à la fois~: le cours de clôture d'une action figure dans le flux de la bourse, dans la base de ticks, dans
l'instantané de la bibliothèque de valorisation, dans les valorisations du service de P\&L et dans chaque tableur qui l'a recopié. Lorsque deux de ces copies
diffèrent, la société a besoin d'une règle qui dit laquelle est juste.

\begin{definition}[Système de référence, source de référence]\label{def:pl:the-platform-map:sor}
Le \emph{système de référence}\index{système de référence} d'un jeu de données est l'unique système dont la copie fait autorité~: les corrections
y sont apportées et nulle part ailleurs, et toute autre copie en est dérivée. Une \emph{source de référence}\index{source de référence} est un système de
référence que la société désigne pour un type de données utilisé par de nombreux systèmes (instruments, contreparties, valorisations de clôture), de sorte que
chaque consommateur lise soit dans ce système, soit dans une copie que la carte montre comme en étant dérivée.
\end{definition}

\begin{proposition}[Un seul propriétaire, une seule réponse]\label{prop:pl:the-platform-map:oneowner}
Si chaque jeu de données possède exactement un \omterm{def:pl:the-platform-map:sor}{système de référence}, si chaque copie est dérivée par la carte de ce système, et si deux rapports lisent la
même version d'un jeu de données, alors ils s'accordent à son sujet.
\end{proposition}

\begin{proof}
La valeur que chaque rapport tire du jeu de données est une fonction de la version lue, appliquée le long des dérivations de la carte~; la même version
par les mêmes dérivations donne la même valeur. Un désaccord exige donc deux \omterm{def:pl:the-platform-map:sor}{systèmes de référence} (deux versions qui
ne dérivent pas l'une de l'autre), une copie faite hors de la carte (une dérivation que personne n'a déclarée), ou deux versions lues à
des instants différents.
\end{proof}

La démonstration est courte parce que la proposition est presque une définition, et son intérêt tient aux trois façons dont elle peut échouer. Les deux
premières sont structurelles et un programme peut les repérer dans la carte~; la troisième concerne le temps, et le reste de ce chapitre lui est consacré.
Le \Cref{lst:pl:the-platform-map:validate} est le contrôle structurel. Sur la carte telle qu'elle a d'abord été dessinée, il relève que deux systèmes
revendiquent les valorisations, la bibliothèque de valorisation et le product control.

\begin{example}[Deux prétentions sur les valorisations]\label{ex:pl:the-platform-map:marks}
La bibliothèque de valorisation calcule les valorisations de fin de journée à partir des cours de clôture et de l'instantané des courbes et des volatilités~; le product
control les vérifie indépendamment par rapport aux cotations des courtiers et aux données de consensus (livre~6, chapitre~27). Les deux appellent
leur sortie «~les valorisations~». Ce sont deux jeux de données~: les valorisations, dont le propriétaire est la bibliothèque de valorisation, et les
valorisations vérifiées avec leurs ajustements, dont le propriétaire est le product control. La correction de la carte consiste à nommer le second jeu de
données et à rendre explicite sa dépendance au premier~; le contrôle passe alors, et un désaccord entre les deux devient un chiffre que le product control
rapporte, et non une découverte.
\end{example}

\omcode{code/firm/platmap/firm_platmap.py}{85}{107}{Contrôles structurels~: un système de référence par jeu de données, noms connus, un
producteur par jeu de données et aucun cycle.}\label{lst:pl:the-platform-map:validate}

\begin{dated}{2026-09}{À quel point les banques sont loin d'un propriétaire unique par donnée}\label{dat:pl:the-platform-map:bcbs}
Les principes du Comité de Bâle pour une agrégation efficace des données de risque et la publication de rapports sur les risques (BCBS~239, publiés le 9 janvier 2013)
demandent aux banques de produire des données de risque exactes, complètes et disponibles en temps voulu, à partir d'une architecture de données qui le permette. Dans son
rapport d'étape du 28 novembre 2023 (consulté en septembre 2026), le Comité a constaté que, sur les 31 banques d'importance systémique mondiale
qu'il a évaluées, seules deux étaient pleinement conformes à l'ensemble des principes.
\end{dated}

\section{Batch et flux}

\begin{definition}[Traitement par lots, traitement en flux, batch de fin de journée]\label{def:pl:the-platform-map:batch}
Le \emph{traitement par lots}\index{traitement par lots} exécute un job sur un jeu de données complet et borné (une journée d'exécutions, un fichier de fournisseur) et
produit sa sortie une seule fois. Le \emph{traitement en flux}\index{traitement en flux} exécute un job en continu sur une suite non bornée
d'événements, en mettant sa sortie à jour à chaque arrivée. Le \emph{batch de fin de journée}\index{batch de fin de journée} est l'ensemble des jobs
batch qui s'exécutent après la clôture pour produire les positions officielles, les valorisations, le P\&L, le risque et les rapports de la journée.
\end{definition}

Une société de trading fait tourner les deux, souvent pour la même grandeur. Le service de positions du chapitre~17 tient les positions en flux pendant
la journée, à partir de chaque exécution au moment où elle arrive, parce que les traders et la barrière de risque en ont besoin tout de suite~; le job de positions de fin de journée recalcule
les positions en batch à partir de la journée complète, du relevé de compensation et du référentiel des titres, parce que les livres sont arrêtés sur ces données.
Les deux doivent concorder, et lorsqu'elles ne concordent pas, l'écart est un écart de rapprochement (chapitre~22).

Dans une chaîne de batchs, les questions portent sur le temps~: à quel instant chaque sortie sera prête, quel job décide de la fin de la chaîne, et avec quel
retard chaque entrée peut-elle arriver. Avec un nombre illimité de machines (le chapitre~13 lève cette hypothèse), deux passages sur la carte y répondent.

\begin{method}[Passes avant et arrière]\label{met:pl:the-platform-map:passes}
Ordonner les jobs de sorte que chaque job vienne après les producteurs de ses entrées.
\emph{Avant~:} pour chaque job, son début est le plus tardif des instants où ses entrées sont prêtes (ou son départ fixe), sa fin est le début plus
la durée, et ses sorties sont prêtes à sa fin~; les entrées externes sont prêtes à leurs instants d'arrivée.
\emph{Arrière~:} pour chaque job, pris dans l'ordre inverse, sa fin au plus tard est la plus petite de sa propre échéance et des débuts au plus tard des
jobs qui lisent ses sorties~; son début au plus tard est cette fin moins sa durée. La \emph{marge} d'un job est son début au plus tard moins son
début planifié. L'arrivée externe au plus tard d'une entrée est le plus petit début au plus tard parmi les jobs qui la lisent.
\end{method}

\begin{proposition}[L'ordonnancement piloté par les dépendances est le plus tôt possible]\label{prop:pl:the-platform-map:earliest}
Avec une capacité illimitée et des durées fixes, le passage avant d'une carte dont tous les jobs démarrent dès que leurs entrées sont prêtes donne à chaque job le
début au plus tôt qu'un ordonnancement quelconque puisse lui donner, et la chaîne se termine à la fin de sa plus longue suite de jobs~: la chaîne qui se termine en dernier,
remontée par l'entrée prête en dernier à chaque étape.
\end{proposition}

\begin{proof}
Par récurrence dans l'ordre des dépendances. Un job ne peut démarrer avant que toutes ses entrées soient prêtes~; si chaque entrée est prête au plus tôt
(hypothèse de récurrence, ou instant d'arrivée pour une entrée externe), le job démarre au plus tôt
en commençant à la plus tardive de ces dates. En remontant depuis le job qui se termine en dernier par l'entrée prête en dernier, on obtient une chaîne de jobs
sans trou entre la fin de l'un et le début du suivant, dont la durée totale est l'instant de fin.
\end{proof}

Le \Cref{lst:pl:the-platform-map:passes} contient les deux passages. Le passage avant enregistre aussi chaque lecture d'une entrée qui n'était pas encore
complète au moment où le job a démarré, ce qui n'est possible que pour un job à départ fixe~: c'est le contrôle dont le crochet avait besoin.

\omcode{code/firm/platmap/firm_platmap.py}{185}{219}{Le passage avant (démarrage sur la dernière entrée, ou à l'heure fixe), les lectures
périmées qu'il implique et le passage arrière des débuts au plus tard.}\label{lst:pl:the-platform-map:passes}

\begin{omfigure}
\begin{tikzpicture}
\begin{axis}[width=12cm,height=8.2cm,xbar stacked,bar width=6pt,xmin=0,xmax=16,
  xtick={0,2,4,6,8,10,12,14,16},xticklabels={16~h~00,18~h~00,20~h~00,22~h~00,00~h~00,02~h~00,04~h~00,06~h~00,08~h~00},
  ytick={0,1,2,3,4,5,6,7,8,9,10,11,12,13,14},yticklabels={backtests,variables,déclarations de transactions,surveillance,validation du P\&L,rapprochement,confirmations,saisie des transactions,batch de risque,scénarios,P\&L,positions,valorisations de clôture,copie de référence,ticks compactés},
  tick label style={font=\scriptsize},ymin=-0.7,ymax=14.7,grid=major,grid style={gray!20},
  legend columns=2,legend style={font=\footnotesize,at={(0.5,-0.12)},anchor=north,draw=none,/tikz/every even column/.append style={column sep=8pt}},
  table/col sep=comma]
\addplot[draw=none,fill=none,forget plot] table[x=start,y=k]{figdata/platforms/01-the-platform-map/gantt_nominal.csv};
\addplot[fill=omDef!45,draw=omDef] table[x=other,y=k]{figdata/platforms/01-the-platform-map/gantt_nominal.csv};
\addplot[fill=omThm!55,draw=omThm] table[x=chain,y=k]{figdata/platforms/01-the-platform-map/gantt_nominal.csv};
\draw[omThm,thick,dashed] (axis cs:15,-0.7) -- (axis cs:15,14.7) node[pos=0.97,left,font=\scriptsize]{07:00 échéance};
\legend{job,chaîne la plus longue}
\end{axis}
\end{tikzpicture}

\omcaption{La chaîne du \omterm{def:pl:the-platform-map:batch}{batch de fin de journée} de la petite société, une nuit nominale, chaque job lancé dès que ses entrées sont prêtes. La chaîne la plus longue (en rouge)
remonte du fichier d'opérations sur titres jusqu'à la copie de référence des données de référence, aux variables de recherche et aux
backtests de nuit, et se termine à 03~h~55~; la ligne en pointillés marque l'échéance de 07~h~00. Données~: \texttt{fig\_platmap.py}.}\label{fig:pl:the-platform-map:gantt}
\end{omfigure}

Une nuit nominale (la \cref{fig:pl:the-platform-map:gantt}) la chaîne se termine à 03~h~55, 185 minutes avant son échéance. Sa chaîne la plus longue
n'est pas le batch de risque, qui est le job le plus long (six heures, de 20~h~15 à 02~h~15), mais la branche recherche~: la copie de référence
attend le fichier d'opérations sur titres de 19~h~30, les variables attendent la copie de référence, et six heures de backtests suivent. Une société qui
aurait acheté davantage de machines pour la grille de risque ne gagnerait pas une minute sur la fin de la nuit. Le \cref{tab:pl:the-platform-map:latest}
transforme la passe arrière en la question que les opérations posent réellement au fournisseur~: à quelle heure au plus tard le fichier peut-il arriver~?

\begin{table}[ht]
\centering\small
\begin{tabular}{lrrr}
\toprule
entrée & arrivée & dernière arrivée sûre & marge (min) \\
\midrule
exécutions, saisie & 16~h~05 & 00~h~40, 22~h~25 & 515, 380 \\
cours de clôture & 16~h~20 & 00~h~40 & 500 \\
courbe et instantané de volatilité & 16~h~45 & 00~h~20 & 455 \\
données de référence du fournisseur & 18~h~00 & 22~h~35 & 275 \\
opérations sur titres du fournisseur & 19~h~30 & 22~h~35 & 185 \\
relevé de compensation & 21~h~30 & 05~h~50 & 500 \\
\bottomrule
\end{tabular}
\caption{Arrivée et dernière arrivée sûre de chaque entrée externe pour chaque rapport du matin devant être prêt à 07~h~00, issues de la
passe arrière sur la carte nominale. Le fichier d'opérations sur titres a la plus petite marge. Données~: \texttt{pl\_platmap.latest\_arrivals}.}\label{tab:pl:the-platform-map:latest}
\end{table}

\section{Lire une carte de plateforme}

Deux autres questions complètent la lecture. La première est structurelle~: ce qu'atteint une entrée~?

\begin{definition}[Rayon d'impact]\label{def:pl:the-platform-map:blast}
Le \emph{rayon d'impact}\index{rayon d'impact} (blast radius) d'un jeu de données est l'ensemble des jeux de données qui en dérivent, directement ou par
l'intermédiaire d'autres jobs~: tout ce qui est faux si lui-même est faux, et tout ce qui est en retard si lui-même est en retard.
\end{definition}

Sur la carte de la petite société, les exécutions de la journée, le fichier de référence du fournisseur et le fichier d'opérations sur titres
atteignent chacun neuf des quinze jeux de données dérivés~; le relevé de compensation en atteint
deux (les écarts de rapprochement et le P\&L signé)~; la capture des données de marché en atteint quatre. Une opération sur titres erronée est aussi dangereuse qu'une exécution erronée, et bien moins surveillée~:
les exécutions sont rapprochées du drop copy de la plateforme (livre~13, chapitre~21) en quelques secondes, alors qu'une opération sur titres
n'est vérifiée, si elle l'est, que lorsqu'un rapport paraît bizarre.

La seconde question porte sur le temps, et c'est celle de l'accroche. Un job qui démarre à une heure fixe lit quelle que soit la version de ses
entrées qui existe à cet instant. Si le producteur d'une entrée est en retard, le job lit la version précédente sans erreur, et sa sortie
est silencieusement en désaccord avec les sorties calculées à partir de la nouvelle.

\begin{proposition}[Quand un démarrage fixe lit des données périmées]\label{prop:pl:the-platform-map:stale}
Un job de démarrage fixe $s$ lit une version complète de l'entrée $x$ si et seulement si $x$ est prête à $s$. Si $x$ est la sortie d'une
chaîne de jobs pilotés par les dépendances, de durée totale $D_x$, qui démarre sur une entrée externe arrivant à l'instant $a$, le job lit
des données périmées exactement lorsque $a > s - D_x$.
\end{proposition}

\begin{proof}
La première partie est la définition de la passe avant. Pour la seconde, la passe avant donne à $x$ l'instant de disponibilité $a + D_x$ lorsque
$a$ est la dernière entrée de la chaîne, et la lecture est périmée exactement lorsque $a + D_x > s$.
\end{proof}

Le batch de risque de l'accroche démarre à 21~h~00 ($s = 300$ minutes après la clôture). Ses entrées comprennent les positions, prêtes 45
minutes après l'arrivée du fichier d'opérations sur titres (25 minutes pour la copie de référence, 20 pour les positions), si bien qu'il lit des positions
périmées dès que le fichier arrive après 20~h~15. Une nuit nominale, le fichier arrive à 19~h~30, avec 45 minutes de marge, et personne ne remarque le
démarrage fixe. La nuit de l'accroche, il est arrivé à 21~h~30~: la copie de référence était prête à 21~h~55, les positions à 22~h~15, et le batch de risque
tournait sur les anciennes versions depuis 21~h~00 (\cref{fig:pl:the-platform-map:hook}). Le job de P\&L, piloté par les dépendances, a lu
les nouvelles à 22~h~15. Le matin des trois équipes a été le prix d'un seul démarrage fixe.

\begin{omfigure}
\begin{tikzpicture}
\begin{axis}[width=12cm,height=4.4cm,xbar stacked,bar width=8pt,xmin=4,xmax=12,
  xtick={4,5,6,7,8,9,10,11,12},xticklabels={20~h~00,21~h~00,22~h~00,23~h~00,00~h~00,01~h~00,02~h~00,03~h~00,04~h~00},
  ytick={0,1,2,3},yticklabels={batch de risque (fixé à 21~h~00),P\&L,positions,copie de référence},
  tick label style={font=\scriptsize},ymin=-0.6,ymax=3.6,grid=major,grid style={gray!20},table/col sep=comma]
\addplot[draw=none,fill=none,forget plot] table[x=start,y=k]{figdata/platforms/01-the-platform-map/gantt_hook.csv};
\addplot[fill=omDef!45,draw=omDef] table[x=dur,y=k]{figdata/platforms/01-the-platform-map/gantt_hook.csv};
\draw[omThm,thick,dashed] (axis cs:5.9167,-0.6) -- (axis cs:5.9167,3.6);
\draw[omThm,thick,dashed] (axis cs:6.25,-0.6) -- (axis cs:6.25,3.6);
\node[font=\scriptsize,omThm,anchor=west] at (axis cs:6.35,1.55) {référentiel des titres prêt à 21\,h\,55, positions à 22\,h\,15};
\end{axis}
\end{tikzpicture}

\omcaption{La nuit de l'accroche~: le fichier d'opérations sur titres arrive à 21~h~30. Le batch de risque, lancé par l'horloge à 21~h~00, a
déjà lu le référentiel des titres et les positions (leurs nouvelles versions sont prêtes aux lignes en pointillés)~; le job de P\&L, lancé dès que ses
entrées sont prêtes, lit les nouvelles. Données~: \texttt{fig\_platmap.py} (\texttt{firm\_platmap.stale\_reads}).}\label{fig:pl:the-platform-map:hook}
\end{omfigure}

Le correctif n'est pas de repousser le démarrage fixe à 22~h~00, ce qui ne fait que déplacer le seuil~; il consiste à faire démarrer le batch de risque
sur ses entrées. Piloté par les dépendances, la nuit de l'accroche, il démarre à 22~h~15 et se termine à 04~h~15, largement dans l'échéance, et la carte
ne signale aucune lecture périmée. Les démarrages fixes existent pour de bonnes raisons (une machine n'est libre qu'après 21~h~00~;
un fichier de fournisseur n'a pas de signal de fin), et lorsqu'on en garde un, le contrôle de la carte transforme le désaccord silencieux en
précondition non satisfaite~: le job refuse de démarrer sur une entrée incomplète, ou démarre et marque sa sortie comme construite sur une version périmée.

\begin{remark}[Ce que la carte ignore]\label{rem:pl:the-platform-map:limits}
Les passes supposent des durées fixes et des machines illimitées. Les durées varient d'une nuit à l'autre, les jobs partagent un cluster, et un job
peut échouer et être relancé~; le chapitre~13 ajoute les files d'attente et les pannes, et le chapitre~20 dimensionne le batch de risque lui-même. La carte ne connaît
aussi que les flux que quelqu'un a déclarés~: un tableur qui recopie les valorisations à 18~h~00 lui est invisible, ce qui explique que la carte soit imposée
en obligeant les systèmes à lire leurs entrées à travers elle (chapitres~6, 7 et~16), plutôt qu'en étant dessinée à côté.
\end{remark}

\section{Tutoriel~: la petite société comme carte}\label{tut:pl:the-platform-map:tut}

\begin{tutorial}
\textbf{Objectif.} Décrire la petite société sous forme de données, la vérifier, planifier une nuit et trouver par programme le défaut de l'accroche.
\textbf{État final~:} \cref{fig:pl:the-platform-map:gantt,fig:pl:the-platform-map:hook}, \cref{tab:pl:the-platform-map:latest}
et le \omterm{def:pl:the-platform-map:blast}{rayon d'impact} de chaque entrée.
\begin{tutsteps}
\item \textbf{Déclarer.} Lire \texttt{pl\_platmap}~: les systèmes par couche, les jeux de données avec leur \omterm{def:pl:the-platform-map:sor}{système de référence} et, pour les entrées externes,
leurs heures d'arrivée, ainsi que les jobs du \cref{lst:pl:the-platform-map:jobs}.
\item \textbf{Valider.} \texttt{draft().validate()} signale que les valorisations ont deux \omterm{def:pl:the-platform-map:sor}{systèmes de référence}~; \texttt{firm().validate()}
ne signale rien.
\item \textbf{Planifier.} \texttt{firm().schedule()} et \texttt{longest\_chain()} donnent la \cref{fig:pl:the-platform-map:gantt}~:
la chaîne copie de référence, variables, backtests, qui se termine à 03~h~55. \texttt{slack(900)} et \texttt{latest\_arrivals} donnent
le \cref{tab:pl:the-platform-map:latest}.
\item \textbf{Rejouer l'accroche.} \texttt{firm(risk\_start=300).stale\_reads(HOOK\_NIGHT)} renvoie les deux lectures périmées du
batch de risque~; \texttt{firm().stale\_reads(HOOK\_NIGHT)} n'en renvoie aucune.
\item \textbf{\omterm{def:pl:the-platform-map:blast}{Rayon d'impact}.} \texttt{firm().blast\_radius()} et \texttt{downstream} du fichier d'opérations sur titres.
\end{tutsteps}
\textbf{Pour aller plus loin.} Donner aux backtests une échéance de 06~h~00 et trouver la nouvelle dernière arrivée du fichier d'opérations sur titres~;
ajouter un job de prévision de marge qui lit les positions et le relevé de compensation et doit être prêt à 06~h~00, puis voir
quelle entrée voit sa dernière arrivée se déplacer.
\end{tutorial}

\section{Construction~: la carte de la plateforme}\label{bld:pl:the-platform-map:platmap}

\begin{build}
\bfield{Rôle} La carte des systèmes, jeux de données et jobs de la société, sous forme de données, avec les contrôles et les questions de temps que chaque
chapitre suivant utilise~: la base de ticks (chapitre~4) et les données de référence (chapitre~6) y déclarent leurs sorties, les règles de qualité des données
du chapitre~7 et l'observabilité du chapitre~28 s'attachent à ses jeux de données, et l'ordonnanceur du chapitre~13 exécute ses jobs.
\bfield{Interface} \texttt{System(name, layer, owner)}, \texttt{Dataset(name, system, cadence, ready)},
\texttt{Job(name, inputs, outputs, duration, start, deadline)}~; \texttt{PlatformMap(systems, datasets, jobs)} avec \texttt{validate},
\texttt{producer}, \texttt{downstream}, \texttt{upstream}, \texttt{schedule}, \texttt{ready\_times}, \texttt{stale\_reads},
\texttt{latest\_starts}, \texttt{slack}, \texttt{longest\_chain}, \texttt{blast\_radius}~; les instants sont exprimés en minutes après la clôture.
\bfield{Règles} Un \omterm{def:pl:the-platform-map:sor}{système de référence} par jeu de données et un producteur par jeu de données dérivé~; aucun cycle~; un job à début fixe est admis
mais toute lecture effectuée avant que son entrée soit prête est signalée~; la carte est le seul endroit où un flux est déclaré.
\bfield{Tests d'acceptation} \texttt{code/firm/platmap/tests/}~: la passe avant et la plus longue chaîne sur une petite carte, et comment une
entrée tardive déplace les deux~; la passe arrière et le slack à la main~; une lecture périmée repérée et absente~; deux \omterm{def:pl:the-platform-map:sor}{systèmes de référence}, un
système inconnu, un jeu de données orphelin et un cycle signalés~; les fermetures et le \omterm{def:pl:the-platform-map:blast}{rayon d'impact}.
\bfield{Pour aller plus loin} Des durées sous forme de lois et la probabilité de respecter l'échéance par simulation~; des versions des jeux de données
et la version lue par chaque job~; une carte générée à partir des déclarations propres aux systèmes plutôt qu'écrite à la main.
\end{build}

\begin{omsources}
\item Basel Committee on Banking Supervision, \emph{Principles for effective risk data aggregation and risk reporting}
(BCBS~239), janvier 2013, et \emph{Progress in adopting the Principles}, novembre 2023.
\item T.~Akidau et al., «~The dataflow model~», \emph{Proceedings of the VLDB Endowment} 8(12), 2015 (le traitement par lots comme cas particulier
du traitement en flux).
\end{omsources}

\section{Exercices}

\begin{exercise}[$\star$]\label{exo:pl:the-platform-map:1}
Le job de P\&L lit les positions (prêtes à 20~h~15), les valorisations (17~h~05) et le référentiel des titres (19~h~55), et dure 25 minutes.
Déterminer à quelle heure il commence et finit une nuit nominale, et quelle entrée le décide.
\end{exercise}

\begin{exercise}[$\star$]\label{exo:pl:the-platform-map:2}
Le fichier de données de référence du fournisseur arrive à 22~h~50 au lieu de 18~h~00. En utilisant le \cref{tab:pl:the-platform-map:latest},
déterminer quel rapport est en retard, et de combien de minutes.
\end{exercise}

\begin{exercise}[$\star$]\label{exo:pl:the-platform-map:3}
Expliquer pourquoi le relevé de compensation a le plus petit \omterm{def:pl:the-platform-map:blast}{rayon d'impact} des sept entrées, et si un petit \omterm{def:pl:the-platform-map:blast}{rayon d'impact}
signifie qu'une entrée demande moins de contrôles.
\end{exercise}

\begin{exercise}[$\star\star$]\label{exo:pl:the-platform-map:4}
Pour finir plus tôt, l'équipe risque avance le début fixe du batch de risque de 21~h~00 à 20~h~30. Déterminer jusqu'à quelle heure
d'arrivée du fichier d'opérations sur titres le batch lit des entrées complètes, et quelle marge cela laisse sur une nuit nominale.
\end{exercise}

\begin{exercise}[$\star\star$]\label{exo:pl:the-platform-map:5}
Les valorisations vérifiées indépendamment par le product control diffèrent de celles de la bibliothèque de pricing de quelques points de base sur certaines positions.
Expliquer pourquoi il s'agit de deux jeux de données et non de deux copies d'un même, et lequel des deux la validation du P\&L doit lire.
\end{exercise}

\begin{exercise}[$\star\star$]\label{exo:pl:the-platform-map:6}
Pendant la journée, le service de positions tient les positions en flux~; après la clôture, un job batch les recalcule. Donner deux raisons
pour lesquelles les deux peuvent différer à 17~h~00, et dire sur laquelle des deux les comptes sont clos.
\end{exercise}

\begin{exercise}[$\star\star\star$]\label{exo:pl:the-platform-map:7}
\emph{Programmation.} Ajouter un job \texttt{margin forecast} qui lit les positions et le relevé de compensation, dure 60 minutes et doit
être prêt à 06~h~00. Calculer son début au plus tard, dire s'il modifie le début au plus tard d'un autre job, et
donner la nouvelle dernière arrivée sûre du relevé de compensation.
\end{exercise}

\begin{exercise}[$\star\star\star$]\label{exo:pl:the-platform-map:8}
\emph{Trouver la faille.} «~Le batch de nuit est trop proche de son échéance. On doublera les machines de la grille de risque, ce qui divise
par deux le batch de risque à trois heures, et tout le batch finira trois heures plus tôt.~»
\end{exercise}

\section{Problème~: Le matin où les chiffres ne concordaient pas}

\begin{problem}[{Problème du week-end --- un fichier en retard et un début fixe}]\label{pb:pl:the-platform-map:1}
La carte de la société miniature (\texttt{pl\_platmap}), sa nuit nominale et la nuit où le fichier d'opérations sur titres est arrivé à 21~h~30.

\medskip\noindent\textbf{Partie I --- La nuit nominale.}
\begin{enumerate}
\item À quelle heure commencent et finissent la copie de référence, les positions et le batch de risque, chaque job démarrant dès que ses entrées sont prêtes~?
\item Quel job finit en dernier, et à quelle heure~?
\item Quelle est la plus longue chaîne, et pourquoi le batch de risque, le job le plus long, n'y figure-t-il pas~?
\item Combien de minutes de slack le batch a-t-il avant 07~h~00~?
\item Quel est le slack du batch de risque~?
\end{enumerate}

\medskip\noindent\textbf{Partie II --- Jusqu'à quelle heure les fichiers peuvent-ils arriver~?}
\begin{enumerate}[resume]
\item Quelle est la dernière arrivée sûre du fichier d'opérations sur titres du fournisseur, et quel job la fixe~?
\item Quelle est la dernière arrivée sûre de la capture des données de marché, et pourquoi est-elle plus tôt que celle des exécutions~?
\item Quelle entrée présente la plus grande marge entre son arrivée nominale et sa dernière arrivée sûre~?
\item Combien de jeux de données le fichier d'opérations sur titres atteint-il~?
\item Quels jeux de données un relevé de compensation erroné atteint-il~?
\end{enumerate}

\medskip\noindent\textbf{Partie III --- La nuit de l'accroche.}
\begin{enumerate}[resume]
\item Le fichier arrive à 21~h~30. Quand le référentiel des titres et les positions sont-ils prêts~?
\item Le batch de risque démarre à 21~h~00 à l'horloge. Lesquelles de ses entrées lit-il à l'état périmé~?
\item Quand tourne le job de P\&L, et quelles versions lit-il~?
\item Jusqu'à quelle heure d'arrivée du fichier le batch de risque à début fixe aurait-il lu des entrées complètes~?
\item Démarré sur ses entrées, quand le batch de risque aurait-il tourné cette nuit-là, et aurait-il respecté l'échéance~?
\end{enumerate}

\medskip\noindent\textbf{Partie IV --- Le verdict.}
\begin{enumerate}[resume]
\item Énoncer le \emph{résultat nommé}~: la fin et le slack de la nuit nominale, la dernière arrivée sûre du fichier
d'opérations sur titres, et l'heure d'arrivée au-delà de laquelle le batch de risque à début fixe lit des données périmées.
\item Pourquoi le désaccord n'a-t-il produit aucune erreur dans aucun système~?
\item Lequel des trois modes de défaillance de la \cref{prop:pl:the-platform-map:oneowner} s'est produit~?
\item Donner deux façons de conserver un début fixe sans le désaccord silencieux.
\item En une phrase~: que donne une carte de la plateforme qu'un schéma de systèmes ne donne pas~?
\end{enumerate}
\end{problem}

\section{Questions d'entretien}

\begin{interviewq}[$\star$ \iqroles{développeur, risque}]\label{iq:pl:the-platform-map:1}
Qu'est-ce qu'un \omterm{def:pl:the-platform-map:sor}{système de référence}, et pourquoi une société en a-t-elle besoin pour chaque jeu de données~?
\end{interviewq}

\begin{interviewq}[$\star$ \iqroles{développeur}]\label{iq:pl:the-platform-map:2}
Dans quels cas calculeriez-vous les positions en flux, et dans quels cas en batch~?
\end{interviewq}

\begin{interviewq}[$\star\star$ \iqroles{développeur, risque}]\label{iq:pl:the-platform-map:3}
Le rapport de risque du matin et le P\&L ne concordent pas sur une position. Comment en trouveriez-vous la cause~?
\end{interviewq}

\begin{interviewq}[$\star\star$ \iqroles{développeur}]\label{iq:pl:the-platform-map:4}
Un job batch doit-il démarrer à heure fixe ou dès que ses entrées sont prêtes~? Qu'est-ce qui pose problème dans chaque cas~?
\end{interviewq}

\begin{interviewq}[$\star\star$ \iqroles{développeur}]\label{iq:pl:the-platform-map:5}
Comment vérifieriez-vous jusqu'à quelle heure un fichier de fournisseur peut arriver sans que les rapports du matin soient en retard~?
\end{interviewq}

\begin{interviewq}[$\star\star\star$ \iqroles{développeur, risque}]\label{iq:pl:the-platform-map:6}
Vous rejoignez une société qui n'a aucun schéma de ses systèmes. Comment construiriez-vous une carte qui reste exacte, et que vérifieriez-vous en premier~?
\end{interviewq}
